\section{Experimental design}
\label{sec:experiments}

\subsection{Controller arms and randomization}

The primary controller comparison (R2-C1) uses a blocked randomized design.
The experimental block is one task instance $\times$ model configuration
$\times$ sealed fault episode with a forkable pre-decision snapshot. Within
each block, all six non-oracle policies are executed from the same snapshot
in randomized order. When exact replay is impossible, independently
generated matched episodes are used with episode included as a random
effect. The primary comparison is the adaptive policy versus the strongest
non-oracle baseline at matched verification budget.

\subsection{Post-training arms}

The primary learning comparison (R2-T) uses six arms, each consuming
different training information from the same pool of adjudicated transitions:

\begin{enumerate}
    \item \textbf{Base model:} no training; establishes unadapted recovery
          and clean-task behavior.
    \item \textbf{Successful-only SFT:} supervised fine-tuning on valid
          terminal trajectories only; tests whether ordinary imitation of
          successful runs is sufficient.
    \item \textbf{Recovery SFT:} supervised fine-tuning on failed and
          verified-repaired transitions; isolates the value of recovery
          examples.
    \item \textbf{Outcome-only optimization:} reinforcement learning with
          independently adjudicated terminal reward; tests whether terminal
          feedback alone is sufficient.
    \item \textbf{Process-supervised optimization:} reinforcement learning
          with independent step-level and transition-level evidence;
          tests the incremental value of process information.
    \item \textbf{Process plus calibrated abstention:} process-supervised
          optimization with post-hoc calibration on a held-out split;
          tests selective safety and over-intervention.
\end{enumerate}

All arms use one open model family with parameter-efficient adaptation
(e.g., Qwen with LoRA). At least three independent training seeds are used
per trainable arm. The number of source trajectories, effective tokens,
optimizer steps, and evaluation budget are matched across arms; unavoidable
mismatches are reported.

\subsection{Matching and checkpoint selection}

Data matching: each trainable arm receives the same number of source
trajectories, drawn from the same development splits. Compute matching:
each arm uses the same number of optimizer steps and the same evaluation
budget. The checkpoint selection rule is fixed before test evaluation
and uses only development-set outcomes. Testing every checkpoint and
selecting the best test result is prohibited.

\subsection{Evaluation sets}

Eight evaluation sets are frozen before training:

\begin{enumerate}
    \item \textbf{Seen task, seen fault:} verify basic learning.
    \item \textbf{Unseen task templates:} test task transfer.
    \item \textbf{Unseen injector implementations:} test mechanism learning
          rather than injector-artifact learning.
    \item \textbf{Unseen fault compositions:} primary recovery
          generalization test.
    \item \textbf{Unknown fault family:} test abstention and safe
          escalation when the fault is truly novel.
    \item \textbf{Clean tasks:} non-inferiority test; the policy should not
          intervene unnecessarily on tasks that would succeed without it.
    \item \textbf{Verifier substitution:} test whether the policy exploits
          a specific verifier implementation rather than recovering
          scientifically.
    \item \textbf{Environment/provider shift:} test robustness to changed
          runtime or model version.
\end{enumerate}

\subsection{Primary estimands and power}

The primary endpoint for H3 is the difference in independently verified
scientific success at a common verification and human-effort budget,
comparing adaptive policy versus the strongest non-oracle baseline.
Supporting endpoints include area under the verified-success versus cost
frontier, irreversible damage rate, false intervention on no-fault tasks,
and oracle regret.

The primary endpoint for H4 is the verified recovery rate on unseen fault
compositions, comparing process-supervised optimization with the stronger
of recovery SFT and outcome-only optimization at matched evaluation cost.
The co-primary safety endpoint is clean-task verified success under a
non-inferiority margin of two percentage points relative to the base model
evaluation.

A $2 \times 2 \times 2$ factorial ablation (R2-A) estimates main effects
and interactions of calibrated R1 belief, active verification, and process
supervision. Where an arm is identical to a primary arm, the exact
checkpoint and result are reused rather than retrained under a new label.

\subsection{Hierarchical analysis and multiplicity}

Uncertainty uses hierarchical bootstrap over task templates and instances,
training seeds, and fault families. At least 2{,}000 bootstrap replicates
are used for primary intervals. H3 and H4 each have one primary comparison;
secondary contrasts use Benjamini-Hochberg FDR within declared families.
Exploratory analyses are labeled and do not inherit confirmatory language.

\subsection{Stopping rules and negative-result fallback}

If process supervision does not improve unseen-fault recovery over the
matched comparator, the negative comparison is reported without redefining
the primary contrast. R2 may still be publishable around the controller
and verified-intervention dataset if H3 and the artifact gates are strong.
If the controller also fails to beat fixed baselines, the intervention
benchmark is released with a documented negative controller result.
